\documentclass[runningheads]{llncs}

\usepackage{stmaryrd}


\newcommand{\abs}[4]{{#1}\, #2\! : \! #3.\, #4}
\newcommand{\absu}[3]{{#1}\, #2.\, #3}
\mathchardef\mhyph="2D % Define a "math hyphen"

% unicode
\usepackage[utf8]{inputenc}

\DeclareUnicodeCharacter{27F6}{\ensuremath{\longrightarrow}}
\DeclareUnicodeCharacter{25B8}{\ensuremath{\blacktriangleright}}
\DeclareUnicodeCharacter{3BC}{\ensuremath{\mu}}
\DeclareUnicodeCharacter{3BB}{\ensuremath{\lambda}}
\DeclareUnicodeCharacter{3B4}{\ensuremath{\delta}}
\DeclareUnicodeCharacter{3C9}{\ensuremath{\omega}}
\DeclareUnicodeCharacter{3B3}{\ensuremath{\gamma}}
\DeclareUnicodeCharacter{3C4}{\ensuremath{\tau}}
\DeclareUnicodeCharacter{21D2}{\ensuremath{\To}}
\DeclareUnicodeCharacter{22B2}{\ensuremath{\lhd}}
\DeclareUnicodeCharacter{2080}{\ensuremath{_0}}
%\definecolor{darkred}{RGB}{100,10,10}

\newcommand{\arr}[2]{#1 \mathbin{\longrightarrow} #2}

\newcommand{\cata}[1]{\llparenthesis #1 \rrparenthesis}

\newcommand{\utp}[0]{\mathcal{U}}
\newcommand{\lin}[0]{\lambda^{\iota\nu}}
\newcommand{\interp}[1]{\llbracket #1 \rrbracket}
\newcommand{\interpl}[0]{\llbracket}
\newcommand{\interpr}[0]{\rrbracket}
\newcommand{\choice}[0]{\xi}
\newcommand{\simto}[0]{\triangleright}
\newcommand{\elcap}[0]{\cap}

\newcommand{\dedu}[0]{\triangleright}
\newcommand{\ctxtup}[1]{\ulcorner #1 \urcorner}
\newcommand{\ctxtdn}[1]{\llcorner #1 \lrcorner}
\newcommand{\ctxtmv}[1]{\langle #1 \rangle}

\newcommand{\pcase}[1]{\noindent\underline{#1:}}

\newcommand{\adbmal}[0]{\reflectbox{$\lambda$}}

\newcommand{\casetri}[0]{\RHD}
\newcommand{\lke}[0]{\lhd}

\newcommand{\lam}[2]{\lambda\, #1.\, #2}
\newcommand{\mal}[2]{\abdmal\, #1.\, #2}
\newcommand{\all}[2]{\forall\, #1.\, #2}
\newcommand{\Lam}[2]{\Lambda\, #1.\, #2}
\newcommand{\thereis}[2]{\exists\, #1.\, #2}
\newcommand{\nuu}[2]{\nu\, #1.\,#2}
\newcommand{\inn}[2]{\langle #1 \rangle_{#2}}
\newcommand{\exel}[2]{\adbmal\,#1.\,#2}
\newcommand{\recc}[2]{\textit{rec}\, #1.\, #2}
\newcommand{\mew}[2]{\mu_{#1}\, #2}
\newcommand{\convl}[0]{\blacktriangleleft}
\newcommand{\convr}[0]{\blacktriangleright}
\newcommand{\conjin}[0]{\smalltriangleleft}
\newcommand{\conjout}[0]{\smalltriangleright}
\newcommand{\dcdot}[0]{\mathbin{\ast}}
\newcommand{\To}[0]{\Rightarrow}
\newcommand{\releq}[0]{\topdoteq}

\newcommand{\omga}[4]{\omega\ #1(#2)\,:\,#3.\ #4}
\newcommand{\clause}[2]{#1 \to #2}

\newcommand{\hysep}[0]{\mathbin{\blacklozenge}}

% \newcommand{\todo}[1]{\textbf{[#1]}}
\newcommand{\rref}[1]{\{\ref{#1}\}}

%\newtheorem{definition}[theorem]{Definition}

% \newenvironment{mycomment}
%     {\color{blue} \it
%     }
%     {
%     }
\newcommand{\shift}[2]{{\uparrow^{#1}_{#2}}}
\newcommand{\textblue}[1]{\textcolor{blue}{#1}}

\newcommand{\lcata}[0]{\ensuremath{\mathcal{L}_{\textit{cata}}}}


\newcommand{\gsep}{\mspace{8mu}\lvert\mspace{8mu}}

% types
\newcommand{\sBool}{\textit{Bool}}
\newcommand{\sNat}{\textit{Nat}}
\newcommand{\sNatList}{\textit{NatList}}
\newcommand{\sAlg}{\Rightarrow}
\newcommand{\sArr}{\to}
\newcommand{\sFunctorApp}[2]{#1 ~ #2} % spaced
% \newcommand{\sFunctorApp}[2]{#1(#2)} % parenthesized
\newcommand{\sMu}[1]{\mu #1}
\newcommand{\sFunctorConSig}[3]{(#1) \sArr \sFunctorApp{#2}{#3}}
\newcommand{\sCoerceSig}[2]{#1 <:: #2}
% terms
\newcommand{\sLet}[3]{\text{let}~ #1 = #2 ~\text{in}~ #3}
% \newcommand{\sCoerce}[2]{\text{coerce}~ #1 ~ #2}
\newcommand{\sCoerce}[2]{#1 - #2}
% coercions
\newcommand{\sCata}[2]{\text{cata}~ #1 ~ #2}
\newcommand{\sRoll}[1]{\text{roll}~ #1}
\newcommand{\sUnroll}[1]{\text{unroll}~ #1}
\newcommand{\sSubCov}[2]{\text{cov}[#1] ~ #2}
\newcommand{\sSubData}[1]{\text{id}[\sMu{#1}]}
\newcommand{\sSubArr}[2]{\text{arr}~ #1 ~ #2}
\newcommand{\sSubAlg}[2]{\text{alg}[#1] ~ #2}
\newcommand{\sSubRecVar}[1]{\text{id}[#1]}
% constructors
\newcommand{\sTrue}{\textit{True}}
\newcommand{\sFalse}{\textit{False}}
\newcommand{\sZero}{\textit{Zero}}
\newcommand{\sSuc}{\textit{Suc}}
\newcommand{\sNil}{\textit{Nil}}
\newcommand{\sCons}{\textit{Cons}}

% helpers
\newcommand{\freshMetavariables}[1]{#1 ~\text{fresh}}
\newcommand{\isDatatype}[1]{#1 ~\text{data}}
\newcommand{\atPos}[1]{^{#1}}
\newcommand{\constructorOf}[4]{(#1 : \sFunctorConSig{#2}{#3}{#4}) \in \mathscr{K}(#3, #4)}
% \newcommand{\atPos}[1]{}

\newcommand{\Match}{\textbf{match }}
\newcommand{\With}{\textbf{ with}}
\newcommand{\Case}{\textbf{case }}
\newcommand{\Func}[1]{\textsc{#1}}

\newcommand{\todo}[1]{{{\contour{red}{\textbf{\color{black}TODO}}} \color{red}#1}}

% \newcommand{\todoHenry}[1]{%
%   \begin{tcolorbox}[
%       colback=olive!15,     % Light olive background
%       colframe=olive!50,    % Muted olive border
%       arc=0mm,            % Sharp corners
%       boxrule=0.5pt,      % Border thickness
%       left=6pt, right=6pt, top=6pt, bottom=6pt, % Inner padding
%       before=\par\smallskip, after=\par\smallskip % Paragraph-like spacing
%     ]
%     \todo{Henry}
%     \\[1em]
%     #1
%   \end{tcolorbox}%
% }
\newcommand{\todoHenry}[1]{%
  \todo{Henry}: #1%
}

% \newcommand{\todoAaron}[1]{%
%   \begin{tcolorbox}[
%       colback=blue!15,     % Light blue background
%       colframe=blue!50,    % Muted blue border
%       arc=0mm,            % Sharp corners
%       boxrule=0.5pt,      % Border thickness
%       left=6pt, right=6pt, top=6pt, bottom=6pt, % Inner padding
%       before=\par\smallskip, after=\par\smallskip % Paragraph-like spacing
%     ]
%     \todo{Aaron}
%     \\[1em]
%     #1
%   \end{tcolorbox}%
% }
\newcommand{\todoAaron}[1]{%
  \todo{Aaron}: #1%
}


\begin{document}

\title{Coercive Subtyping via Logic Programming for Painless Mendler-style Algebraic Programming}

%% \author{Henry Blanchette}
%% \email{blancheh@umd.edu}
%% \affiliation{%
%%   \institution{The University of Maryland}
%%   \city{College Park}
%%   \state{Maryland}
%%   \country{USA}
%% }

%% \author{Aaron Stump}
%% \email{stumpaa@bc.edu}
%% \affiliation{%
%%   \institution{Boston College}
%%   \city{Boston}
%%   \state{Massachusetts}
%%   \country{USA}
%% }

%% \begin{CCSXML}
%% <ccs2012>
%% <concept>
%% <concept_id>10011007.10011006.10011008.10011009.10011012</concept_id>
%% <concept_desc>Software and its engineering~Functional languages</concept_desc>
%% <concept_significance>500</concept_significance>
%% </concept>
%% <concept>
%% <concept_id>10011007.10011006.10011008.10011024.10011033</concept_id>
%% <concept_desc>Software and its engineering~Recursion</concept_desc>
%% <concept_significance>500</concept_significance>
%% </concept>
%% </ccs2012>
%% \end{CCSXML}

%% \ccsdesc[500]{Software and its engineering~Functional languages}
%% \ccsdesc[300]{Software and its engineering~Recursion}
%%
%% \keywords{strong functional programming, type-based termination, Mendler-style recursion}


\maketitle

\begin{abstract}
  The paper proposes coercive subtyping to infer combinators needed
  for algebraic programming.  

\keywords{Strong functional programming, Mendler-style recursion, coercive subtyping}
\end{abstract}



\section{Introduction} % Aaron

% Algebraic programming as an example of total functional programming (cite Turner)
% Mendler-style algebra
% Burden of inserting roll/unroll and cata
% We will reduce this burden using coercive subtyping to infer the coercions

Strong functional programming is a form of pure functional programming
where all functions are statically required to terminate on all
inputs.  One way to achieve this is through an algebraic approach:
datatypes are least fixed points $\mu F$ of signature functors $F$,
and recursive functions are least fixed points $\cata{a}$ of algebras
$a$.  The functions are called catamorphisms.  The signature functor
can be seen as describing one layer of the data.  The least fixed
point then describes data consisting of any finite number of layers.
We may unroll $\mu F$ to $F\ (\mu F)$, and roll it back again.  Such
transformations are needed to allow construction of data by applying
constructors of the signature functor, and destruction by
pattern-matching against those constructors.  Algebras interpret the
operations declared for one layer of the data, and catamorphisms then
apply algebras across all the layers.  There is no other mechanism for
recursion or looping in the language.  All functions written in this
style are guaranteed to terminate.

One hindrance to algebraic programming is the need to apply functions
for rolling and unrolling signature functors, and to apply $\cata{-}$
to obtain a recursive function from an algebra.


\section{Motivating example} % Henry

% pick nice example
%
%   subtraction uses unroll, but not roll or cata
%   length uses roll, but not unroll or cata
%   write an outer function that subtracts some quantity from the length of a list
%
% here it is in Haskell
% here it is in algebraic style with coercions
%
% Can we use the algebraic style with lighter syntax?  This would make it more
% feasible in practice.

\section{Syntax} % Henry (maybe Aaron helping with prose)

\section{Constraint-Generating Typing} % Henry

\section{Subtyping} % Aaron

  \subsection{Declarative}

  \subsection{Algorithmic}

  \subsection{Equivalence of declarative and algorithmic subtyping}

\section{Implementation with Chronolog} % Henry

\section{Returning to the Examples} % Henry

% remind the version with coercions
% show without coercions

\section{Related Work} % Aaron

Turner introduced the term \emph{strong functional programming},
presented several arguments in favor of the approach, and proposed an
\emph{elementary} way to realize such a language, in the form of
static checks for structural recursion~\cite{turner95}.  Mendler
proposed a recursor using an abstract type to ensure that recursive
calls are permitted only on subdata for an inductive
type~\cite{mendler91}.

\section{Conclusion and future work} % Aaron

% Richer patterns of recursion, as in POPL 2023 paper

\bibliographystyle{plain}
\bibliography{paper.bib}

\end{document}
